\section{Podstawy rachunku wektorowego}

\subsubsection*{Skalar}

\begin{equation}
    a
\end{equation}

\subsubsection*{Wektor}

\begin{equation}
    \Vec{a} = [ a_x;\ a_y;\ a_z ] = \begin{bmatrix} a_x \\ a_y \\ a_z \end{bmatrix}
\end{equation}

\subsubsection*{Punkt}

\begin{equation}
    O( O_x;\\ O_y;\\ O_z )
\end{equation}

\subsubsection*{Suma wektorów}

\begin{equation}
    \Vec{a} + \Vec{b} = [ a_x + b_x;\ a_y + b_y;\ a_z + b_z ]
\end{equation}

\subsubsection*{Iloczyn skalarny}

\begin{equation}
    \Vec{a} \cdot \Vec{b} = a_x b_x + a_y b_y + a_z b_z
\end{equation}

\begin{equation}
    \Vec{a} \cdot \Vec{b} = |\Vec{a}| |\Vec{b}| \cos\alpha
\end{equation}

\subsubsection*{Moduł wektora}

\begin{equation}
    |\Vec{a}| = \sqrt{\Vec{a} \cdot \Vec{a}} = \sqrt{ a_x^2 + a_y^2 + a_z^2 }
\end{equation}

\subsubsection*{Wersor}

\begin{equation}
    \Vec{e} = [ e_x;\ e_y;\ e_z ]
\end{equation}

\begin{equation}
    |\Vec{e}| = \sqrt{\Vec{e} \cdot \Vec{e}} = \sqrt{ e_x^2 + e_y^2 + e_z^2 } = 1
\end{equation}

\subsubsection*{Wersor bazowy}

\begin{align}
    \Vec{i} = [1;\ 0;\ 0] \nonumber \\
    \Vec{j} = [0;\ 1;\ 0] \nonumber \\
    \Vec{k} = [0;\ 0;\ 1]
\end{align}

\subsubsection*{Inne zapisy wektora}

\begin{equation}
    \Vec{a} = a_x \Vec{i} + a_y \Vec{j} + a_z \Vec{k}
\end{equation}

\begin{equation}
    \Vec{a} = a \Vec{e}_a = a [ e_{a,x};\ e_{a,y};\ e_{a,z} ]
\end{equation}